% !TEX root = ../main.tex

\chapter[{闭链的上同调类 [Cycle]}]{闭链的上同调类}\label{IV}
\blfootnote{A.\ Grothendieck 作, P.\ Deligne 整理}

本讲受到 Grothendieck 的笔记的启发, 这些笔记构成了 \cite[IV]{sga5} 的初稿.
我们定义域上有限型分离光滑概形中闭链 (cycle) $X$ 的上同调类,
并证明相交对应于杯积 (cup product).

第 \ref{IV:1} 节包含一些一般性的形式准备. 第 \ref{IV:2} 节在若干比上述情形更一般的情况下定义闭链的类.
这里没有讨论的主要相容性, 是闭链的正像与上同调中的迹态射之间的相容性.
第 \ref{IV:3} 节运用这一形式体系, 推出代数闭域 $k$ 上紧合光滑概形的、具有孤立不动点的自同态的 Lefschetz 迹公式,
以及曲线的 Frobenius 自同态的迹公式.

我们采用下列约定:
\begin{enumerate}[\indent 1)]
  \item “概形”指分离 Noether 概形 (这主要是为方便而作的假设).
  \item 第 \ref{IV:2} 节固定整数 $n$, 并假定 $n$ 在所有考虑的概形上均可逆.
  \item 第 \ref{IV:3} 节固定素数 $\ell$, 并假定 $\ell$ 在所有考虑的概形上均可逆.
    所用上同调始终为 $\ell$ 进上同调:
    \[
      \h^\bullet(X) = \dQ_\ell\otimes_{\dZ_\ell} \varprojlim \h^\bullet(X,\dZ/\ell^n) \text{.}
    \]
\end{enumerate}
闭链的上同调类、迹态射等, 都由有限系数 $\dZ/\ell^n$ 的情形取极限定义 (参见 \cite[VI]{sga5}).

\section{带支集上同调与杯积}\label{IV:1}

本节回顾一般拓扑学中的一些内容, 读者只需在用到时查阅.

\subsection{\texorpdfstring{$H^1$}{H1} 与挠子}\label{IV:1-1}

\subsubsection{}\label{IV:1-1-1}

设 $\sF$ 为景 $X$ 上的阿贝尔层. 已知 $\h^1(X,\sF)$ 分类 $X$ 上的 $\sF$-挠子 (torsor).
我们将同构 $($ $\sF$-挠子的同构类之集 $)\to \h^1(X,\sF)$ 规范化 ($=$ 选定其符号),
使得对于每个正合列 $0\to \sF\xrightarrow\alpha\sG\xrightarrow\beta\sH\to 0$
以及每个 $h\in\h^0(X,\sH)$, $\sF$-挠子 $\beta^{-1}(h)\subset \sG$ 的类为 $\partial h$;
其中 $\sF$ 的作用由 $(f,x)\mapsto \alpha(f)+x$ 给出.

\subsubsection{}\label{IV:1-1-2}

设 $\sP$ 为一个 $\sF$-挠子. 若 $(U_i)$ 是 $X$ 的开覆盖, $p_i$ 是 $\sP$ 在 $U_i$ 上的截面,
则给 $\sP$ 配以 \v{C}ech 上闭链
\[
  p_{i j} = p_j - p_i \qquad \text{($p_{i j}\in \h^0(U_i\times U_j,\sF)$).}
\]
若按通常规则定义映射
$\operatorname{\check H}^\bullet(X,\sF)\to \h^\bullet(X,\sF)$,
使之成为 $\delta$-函子的态射, 则
$(p_{i j})\in \operatorname{\check H}^1(X,\sF)$ 在 $\h^1(X,\sF)$ 中的像,
就是按 \ref{IV:1-1-1} 规范化的 $\sP$ 的类.

\subsubsection{}\label{IV:1-1-3}

$\h^\bullet(X,\sF)$ 的定义如下: 若 $\sF^\bullet$ 是 $\sF$ 的一个各项均无高阶上同调的消解,
则 $\h^i(X,\sF) = \h^i\Gamma(X,\sF^\bullet)$.
$\delta$-函子结构由以下方式得到: 给层的短正合列配以消解的短正合列,
且该列在施加函子 $\Gamma$ 后仍正合. 若 $\sF^\bullet$ 是 $\sF$ 的消解,
则与
\[\xymatrix{
  0 \ar[r]
    & \sF \ar[r]
    & \sF^0 \ar[r]^-d
    & \ker(d) \ar[r]
    & 0
}\]
相联系的连接同态 $\partial$ 诱导定义同构
$\h^1(X,\sF) = \Gamma\left(X,\ker(d)\right)/d \Gamma(X,\sF^0)$ 的负值.

\subsubsection{}\label{IV:1-1-4}

设 $U$ 是 $X$ 的开部分 (终层的一个子层), $D$ 为其“闭补集”.
已知 $\h_D^1(X,\sF)$ 分类 $X$ 上在 $U$ 上给定平凡化的 $\sF$-挠子.
对于每个短正合列 $0\to \sF\xrightarrow\alpha \sG\xrightarrow\beta\sH\to 0$
以及每个支集在 $D$ 中的截面 $h\in\h_D^0(X,\sH)$,
由截面 $0$ 在 $U$ 上平凡化的挠子 $\beta^{-1}(h)$ 的类为 $\partial h$.

带支集上同调的长正合列
\[\xymatrix{
  \cdots \ar[r]^-\partial
    & \h_D^i(X,\sF) \ar[r]
    & \h^i(X,\sF) \ar[r]
    & \h^i(U,\sF) \ar[r]^-\partial
    & \cdots
}\]
由函子列
\[\xymatrix{
  0 \ar[r]
    & \Gamma_D \ar[r]
    & \Gamma \ar[r]
    & \Gamma(U,-) \ar[r]
    & 0
}\]
定义 (它在内射层上正合). 对每个截面 $f\in \h^0(U,\sF)$,
$\partial f\in \h_D^1(X,\sF)$ 是由截面 $f$ 在 $U$ 上平凡化的平凡挠子 $\sF$ 的类.

\subsubsection{}\label{IV:1-1-5}

设 $j:U\hookrightarrow X$. 若 $\sF$ 嵌入 $j_* j^* \sF$, 则正合列
\[\xymatrix{
  0 \ar[r]
  & \sF \ar[r]
  & j_* j^* \sF \ar[r]
  & j_* j^* \sF / \sF \ar[r]
  & 0
}\]
给出
$\partial:\h^0(j_* j^*\sF/\sF) = \h_D^0(j_* j^* \sF/\sF) \to \h_D^1(X,\sF)$.
复合映射
\[\xymatrix{
  \h^0(U,\sF) = \h^0(X,j_* j^*\sF) \ar[r]
    & \h^0(X,j_* j^*\sF / \sF) \ar[r]^-\partial
    & \h_D^1(X,\sF)
}\]
是 \ref{IV:1-1-4} 中所考虑映射的负值.

\subsubsection{}\label{IV:1-1-6}

回顾: 若 $\sL$ 是概形 $X$ 上的可逆层, 则与其对应的 $\dG_m$-挠子是层
$\operatorname{Isom}(\sO,\sL)$, $\dG_m$ 通过
$(\lambda,f)\mapsto f\circ (\lambda\cdot) = \lambda f$ 作用于它.
还要回顾: 若 $D$ 是 $X$ 上的 Cartier 除子, $j:U\hookrightarrow X$ 是补开子概形的包含映射,
则可逆层 $\sO(D)$ 是 $j_*\sO_U$ 的子层, 由满足 $s f$ 属于 $\sO_X$ 的局部截面 $s$ 组成,
其中 $f$ 是 $D$ 的局部方程.

\subsection{杯积}\label{IV:1-2}

本小节讨论带支集上同调中杯积的一些性质, 在 [\nameref{V}] 中还会用到它们,
以联系曲线的 Poincaré 对偶性与 Jacobian 的自对偶性.

\subsubsection{}\label{IV:1-2-1}

设 $X$ 为景, $Y$ 为 $X$ 的闭部分, $\sF,\sG$ 为 $X$ 上的两个阿贝尔层.
例如: $X$ 为概形, $Y$ 为闭子概形, $\sF,\sG$ 为 $\et X$ 上的层. 定义乘积
\begin{equation*}\tag{1.2.1.1}\label{IV:eq:1-2-1-1}
  \Gamma_Y(X,\sF)\otimes \Gamma(Y,\sG) \to \Gamma_Y(X,\sF\otimes\sG) \text{.}
\end{equation*}
支集在 $Y$ 中的 $\sF$ 的截面 $s$, 与 $\sG$ 在 $Y$ 上的截面 $t$ 的乘积如下构造:
在 $X$ 上局部地, $t$ 是 $\sG$ 在 $X$ 上的某截面 $t'$ 到 $Y$ 的限制, 于是取乘积 $s\otimes t'$.
其支集在 $Y$ 中, 且不依赖于 $t'$ 的选取; 这说明定义合理, 并使之能够整体化.

设 $i:Y\hookrightarrow X$ 为包含态射. \eqref{IV:eq:1-2-1-1} 的局部对应物为乘积
\begin{equation*}\tag{1.2.1.2}\label{IV:eq:1-2-1-2}
  i^!\sF\otimes i^*\sG \to i^!(\sF\otimes\sG) \text{,}
\end{equation*}
对其施加 $\Gamma(Y,-)$ 就得到 \eqref{IV:eq:1-2-1-1}.

\subsubsection{}\label{IV:1-2-2}

现在将这些箭头导出. 设 $\sK$, $\sL$, $\sM$ 为导出范畴中的对象,
并给定双线性映射 $\sK\lotimes\sL\to\sM$. 将 \eqref{IV:eq:1-2-1-1} 导出, 得到
\begin{equation*}\tag{1.2.2.1}\label{IV:eq:1-2-2-1}
  \eR\Gamma_Y(X,\sK) \lotimes \eR\Gamma(Y,\sL) \to \eR\Gamma_Y(X,\sM)
\end{equation*}
它诱导
\begin{equation*}\tag{1.2.2.2}\label{IV:eq:1-2-2-2}
  \h_Y^i(X,\sK)\otimes \h^j(Y,\sL) \to \h_Y^{i+j}(X,\sM)
\end{equation*}
(杯积). 局部箭头 \eqref{IV:eq:1-2-1-2} 给出
\begin{equation*}\tag{1.2.2.3}\label{IV:eq:1-2-2-3}
  \eR i^! \sK\lotimes i^*\sL \to \eR i^!\sM \text{,}
\end{equation*}
对其施加 $\eR\Gamma(Y,-)$ 就得到 \eqref{IV:eq:1-2-2-1}.

\subsubsection{}\label{IV:1-2-3}

上面, 我没有提及在同一公式中使用右导出函子 ($\eR\Gamma$) 与左导出函子 ($\lotimes$) 所产生的冲突.
\begin{enumerate}[\indent a)]
  \item 下一小节中, 所考虑的右导出函子均具有有限上同调维数.
    这使我们能够始终在导出范畴 $\eD^-$ 中工作.
    为得到既平坦又松弛的复形, 使用 \cite[XVII]{sga4} 中的典范松弛消解.
  \item 要建立更一般的理论, 就不能再将由 $\sK$ 和 $\sL$ 到 $\sM$ 的双线性映射解释为由
    $\sK\lotimes\sL$ 到 $\sM$ 的态射.
    例如, 若两个因子都在 $\eD^+$ 中, 则
    $\eR\Gamma_Y(X,\sK)\lotimes\eR\Gamma(Y,\sL)$ 没有定义.
    一种解决办法是在 $\eD^+$ 中工作, 定义
    \[
      \operatorname{Bil}(\sK,\sL;\sM) = \varinjlim \hom(\sK'\otimes\sL',\sM') \text{,}
    \]
    其中极限取遍拟同构 $\sK'\iso \sK$, $\sL'\iso \sL$, $\sM\iso \sM'$
    ($\sK'$, $\sL'$, $\sM'$ 下有界), $\hom$ 表示“模去同伦的复形态射”,
    并始终使用这种双线性映射, 而完全不提 $\otimes$.
\end{enumerate}

\subsubsection{第二主题}\label{IV:1-2-4}

设 $U$ 为 $X$ 的开子概形, $j:U\hookrightarrow X$ 为包含态射.
对 $U$ 上导出范畴中的 $\sK$, 置
$\eR\Gamma_!(U,\sK)=\eR\Gamma(X,j_!\sK)$.
设 $\sK$, $\sL$, $\sM$ 在 $U$ 上, 并有双线性映射 $\sK\lotimes\sL\to\sM$, 我们希望定义
\begin{equation*}\tag{1.2.4.1}\label{IV:eq:1-2-4-1}
  \eR\Gamma(U,\sK)\lotimes\eR\Gamma_!(U,\sL)\to\eR\Gamma_!(U,\sM) \text{.}
\end{equation*}
在层及其整体截面的层次上, 这样的乘积来自同构
$j_*\sF\otimes j_!\sG\xleftarrow\sim j_!(\sF\otimes \sG)$;
但须注意, $\eR\Gamma_!$ 一般不是 $\Gamma(X,j_!-)$ 的导出函子.

我们先定义
\[\xymatrix{
  \eR j_* \sK\lotimes j_! \sL
    & \ar[l]_-\sim j_!(\sK\lotimes \sL) \ar[r]
    & j_!\sM \text{.}
}\]
施加 $\eR\Gamma(X,-)$, 得到
\[\xymatrix{
  \eR\Gamma(U,\sK)\lotimes\eR\Gamma_!(U,\sL) = \eR\Gamma(X,\eR j_*\sK)\lotimes \eR\Gamma(X,j_!\sL) \ar[r]
    & \eR\Gamma(X,j_! \sM) \text{.}
}\]

这里同样将在下一小节通过在 $\eD^-$ 中工作, 来解决左、右之间的冲突.

\subsubsection{尾声}\label{IV:1-2-5}

设 $j:U\hookrightarrow X$ 为开部分, $i:Y\hookrightarrow U$ 为 $U$ 的闭部分.
设 $\bar Y$ 为 $X$ 的闭子集, 满足 $\bar Y\cap U=Y$ (例如 $U\setminus Y$ 的补集).
设 $\sK$ 在 $U$ 上. 置
$\eR\Gamma_{Y!}(U,\sK)=\eR\Gamma_!(Y,\eR i^!\sK)$,
其中 $\eR\Gamma_!$ 是相对于 $Y$ 到 $\bar Y$ 的包含而言的.
对于 $U$ 中包含 $Y$ 的每个开子集 $V$, $\eR\Gamma_{Y!}(U,\sK)$ 映到 $\eR\Gamma_!(V,\sK)$:
若仍以 $i$ 表示 $\bar Y$ 到 $X$ 的包含, $j$ 表示 $Y$ 到 $\bar Y$ 的包含,
$k$ 表示 $V$ 到 $X$ 的包含, 则有
$\eR i^!\sK = j^*\eR i^! k_!(k^*\sK)$, 从而有态射
$j_!\eR i^!\sK \to \eR i^! k_!(k^*\sK)$.
对其施加 $\eR\Gamma(\bar Y,-)$, 得到
$\eR\Gamma_{Y!}(U,\sK)\to \eR\Gamma_{\bar Y}(k_! k^* \sK) \to \eR\Gamma(k_! k^*\sK) = \eR\Gamma_!(V,\sK)$.

对于 $U$ 上的 $\sK,\sL,\sM$ 及双线性映射
$\sK\lotimes\sL\to\sM$, 我们希望定义
% 据法文原书印页135 (PDF140), 首因子带 Y 支集; 补回英文转录遗漏的下标.
\begin{equation*}\tag{1.2.5.1}\label{IV:eq:1-2-5-1}
  \h_Y^n(U,\sK)\otimes \h_!^m(Y,\sL) \to \h_{Y!}^{n+m}(U,\sM)
\end{equation*}
(最后这个群本身又映到 $\h_!^{n+m}(V,\sM)$).
在导出范畴中, 这意味着要定义
\begin{equation*}\tag{1.2.5.2}\label{IV:eq:1-2-5-2}
  \eR\Gamma_Y(U,\sK) \lotimes \eR\Gamma_!(Y,\sL) \to \eR\Gamma_{Y!}(U,\sM) \text{.}
\end{equation*}
我们将 $\eR\Gamma_Y$ 与 $\eR\Gamma(Y,\eR i^!\sK)$ 等同.
所求的乘积于是具有 \eqref{IV:eq:1-2-4-1} 的形式, 对应于 $Y$ 上的局部乘积
\eqref{IV:eq:1-2-2-3}:
$\eR i^! \sK\lotimes\sL \to \eR i^! \sM$.

\subsubsection{}\label{IV:1-2-6}

上面建立了绝对形式体系. 还有平行的相对形式体系:
将 $\Gamma$ 换成 $f_*$ 即可, 其中 $f$ 是态射 $X\to S$.

\subsection{Koszul 规则}\label{IV:1-3}

设 $A$ 为交换环, $(V_i)_{i\in I}$ 为有限族分次 (或 $\dZ/2$-分次) $A$-模.
回顾 Koszul 规则意义下的分次张量积 $\bigotimes_{i\in I} V_i$ 的定义
(参见 \cite[XVII 1.1]{sga4}). 对 $I$ 上的每个全序 $a$, 我们定义一个模 $V(a)$.
还要定义同构的传递系统 $\varphi_{a b}:V(b)\iso V(a)$;
各 $V_i$ 的分次张量积就是这一模系统的“共同值” $\varprojlim V(a)$. 取
\begin{enumerate}[\indent a)]
  \item $V(a) = \bigotimes_{i\in I} |V_i|$ (各 $V_i$ 的底层非分次模的通常张量积).
  \item 若 $x_i\in V_i$ 为齐次元, 则取
    $\varphi_{a b}(\otimes x_i) = (-1)^N \otimes x_i$,
    其中 $N$ 是对满足 $i<_a j$ 且 $i>_b j$ 的各对 $(i,j)$ 求和所得的 $\deg(x_i)\deg(x_j)$ 之和.
\end{enumerate}

\subsubsection*{\hypertarget{sga4half.e04.example.1-3-1}{例 1}}

取 $I=\{1,2\}$, $a$ 为 $1<2$ 的序, $b$ 为 $1>2$ 的序.
设 $v_i\in V_i$ 为齐次元, 以 $v_1\otimes v_2$ (分别以 $v_2\otimes v_1$)
表示各 $v_i$ 在 $V(a)$ (分别在 $V(b)$) 中乘积在分次张量积中的像. 有
\[
  v_1\otimes v_2 = (-1)^{\deg(v_1)\deg(v_2)} v_2\otimes v_1
\]
(Koszul 规则).

\subsubsection*{\hypertarget{sga4half.e04.example.1-3-2}{例 2}}

若所有 $V_i$ 的次数均为 $1$, 则分次张量积的底层模与各 $V_i$ 的底层模 $|V_i|$ 的通常张量积之间有典范同构:
\[
  \left|\bigotimes V_i\right| \simeq \bigotimes |V_i|\otimes_{\dZ}\bigwedge^{|I|}\dZ^I \text{.}
\]

\subsubsection*{\hypertarget{sga4half.e04.example.1-3-3}{例 3}}

若 $A$ 为域, $X_i$ 为有限族空间, 则 K\"unneth 公式写为
$\h^\bullet\left(\prod X_i,A\right) = \bigotimes \h^\bullet(X_i,A)$.

设 $V^\vee$ 为 $V$ 的分次对偶. 典范映射
\begin{equation*}\tag{1.3.1}\label{IV:eq:1-3-1}
  V^\vee\otimes V = V\otimes V^\vee \to A
\end{equation*}
为 $v'\otimes v\to v'(v)$.

现在假定 $A$ 为域, 并且只考虑有限维向量空间. 典范同构
\begin{equation*}\tag{1.3.2}\label{IV:eq:1-3-2}
  W\otimes V^\vee \to \hom(V,W)
\end{equation*}
为 $w\otimes v'\mapsto \left(v\mapsto w\cdot v'(v)\right)$.
通过此同构, 复合
$\hom(Y,Z)\otimes \hom(X,Y)\to \hom(X,Z)$
等同于由 \eqref{IV:eq:1-3-1} 诱导的态射:
$Z\otimes Y^\vee\otimes Y\otimes X^\vee \to Z\otimes X^\vee$.

$f:V\to V$ 的迹 (齐次 $f$ 的次数为 $\ne 0$ 时迹为零) 是 $f$ 在
\begin{equation*}\tag{1.3.4}\label{IV:eq:1-3-4}
  \tr:\hom(V,V) {\xleftarrow\sim}_\text{\eqref{IV:eq:1-3-2}} V\otimes V^\vee = V^\vee \otimes V \to_\text{\eqref{IV:eq:1-3-1}} A .
\end{equation*}
下的像.

容易验证, 当 $f$ 的次数为 $0$ 时,
\begin{equation*}\tag{1.3.5}\label{IV:eq:1-3-5}
  \tr(f,V) = \sum (-1)^i \tr(f,V^i) \text{.}
\end{equation*}

将由态射 \eqref{IV:eq:1-3-1}
$V^\vee\otimes V\otimes W^\vee\otimes W\to k$
复合得到的两个态射的交换性写出, 可得对齐次映射 $f:V\to W$ 和 $g:W\to V$, 有
\begin{equation*}\tag{1.3.6}\label{IV:eq:1-3-6}
  \tr(f g) = (-1)^{\deg(f)\deg(g)} \tr(g f) \text{.}
\end{equation*}

\section{闭链的上同调类}\label{IV:2}

\subsection{除子的类}\label{IV:2-1}

\subsubsection{}\label{IV:2-1-1}

设 $D$ 为概形 $X$ 中的 Cartier 除子. 在 $D$ 之外, 可逆层 $\sO(D)$ 由截面 $1$ 平凡化.
$D$ 在 $\h_D^1(X,\dG_m)$ 中的类 $\cl(D)$,
是相应的、在 $X\setminus D$ 上给定平凡化的 $\dG_m$-挠子的类
(\ref{IV:1-1-6} 和 \ref{IV:1-1-4}).

% 据法文原书印页138 (PDF143), 边界映射的靶次数为 i+1; 修复英文转录错误.
设 $\partial:\h^i(X\setminus D,\dG_m)\to \h_D^{i+1}(X,\dG_m)$ 为 \ref{IV:1-1-4} 的态射.
若 $D$ 有整体方程 $f$, 则乘以 $f$ 给出从在 $X\setminus D$ 上由 $1$ 平凡化的 $\sO(D)$,
到在 $X\setminus D$ 上由 $f$ 平凡化的 $\sO$ 的同构.
因此由 \ref{IV:1-1-4} 得
\begin{equation*}\tag{2.1.1}\label{IV:eq:2-1-1}
  \cl(D) = \partial f \text{.}
\end{equation*}

对于使 $u^* D$ 仍为 Cartier 除子的每个态射 $u:X'\to X$
(即 $u^{-1} D$ 与 $\operatorname{Ass}(X')$ 不相交), 有
$\cl(u^* D) = u^*\cl(D)$.
若希望对每个态射 $u$ 都有这样的函子性, 则应考虑更一般的带截面的可逆层, 而不只考虑 Cartier 除子.

回顾, 从现在起假定整数 $n$ 在所考虑的概形上可逆.
设 $\partial:\h_D^i(X,\dG_m) \to \h_D^{i+1}(X,\dmu_n)$
为 Kummer 正合列 $0 \to \dmu_n \to \dG_m \to \dG_m \to 0$ 的边界映射.

\begin{definition}\label{IV:2-1-2}
$D$ 在 $\h_D^2(X,\dmu_n)$ 中的\emph{类 $\cl_n(D)$} 为 $\partial \cl(D)$.
\end{definition}

在不致混淆时省略下标 $n$.

\subsubsection{}\label{IV:2-1-3}

图
\[\xymatrix{
  \h^0(X\setminus D,\dG_m) \ar[r]^-\partial \ar[d]^-\partial
    & \h_D^1(X,\dG_m) \ar[d]^-\partial \\
  \h^1(X\setminus D,\dmu_n) \ar[r]^-\partial
    & \h_D^2(X,\dmu_n)
}\]
是反交换的. 因此, 若 $D$ 有整体方程 $f$, 则 $\cl_n(D)$
是 $f$ 的 $n$ 次根所成的 $\dmu_n$-挠子在 $\h^1(X\setminus D,\dmu_n)$ 中的类经 $\partial$ 映射所得像的负值.

\begin{proposition}\label{IV:2-1-4}
设 $i$ 为 $D$ 到 $X$ 的包含. 若 $D$ 和 $X$ 都正则,
则带支集上同调层 $\eR^p i^! \dmu_n$ 在 $p=0,1$ 时为零,
且 $\eR^2 i^!\dmu_n = \underline{\dZ/n}$, 由 $\cl_n(D)$ 生成.
\end{proposition}

只需证明, 当 $X$ 严格局部且 $D$ 由一个正则参数定义时,
$\h_D^p(X,\dmu_n) = 0$ 对 $p=0,1$ 成立,
而 $\h_D^2(X,\dmu_n) = \dZ/n$, 由 $\cl(D)$ 生成.
以 $\sim$ 表示约化上同调, 则有
$\widetilde\h^{p-1}(X\setminus D,\dmu_n)\iso \h_D^p(X,\dmu_n)$.
$p=0,1$ 的断言是说 $D$ 不使 $X$ 不连通; $p=2$ 时, 由 Abhyankar 引理结合 \ref{IV:2-1-3} 得到.

这是相对定理 (\hyperref[I]{平展上同调}, \ref{I:5-3-4}) 的部分对应物.
Grothendieck 猜想 $\eR^p i^! \dmu_n$ 在 $p\ne 2$ 时消失,
至少对优良的 $X$ 如此 (纯性猜想); 但这只在特征 $0$ 的情形已知 (\cite[XIX]{sga4}).

\begin{theorem}[基本相容性]\label{IV:2-1-5}
设 $X$ 为代数闭域 $k$ 上的光滑曲线, $P$ 为 $X$ 的闭点,
$\tr$ 为复合 $\h_P^2(X,\dmu_n) \to \h_c^2(X,\dmu_n) \xrightarrow{\tr} \dZ/n$. 则有
\[
  \tr \cl(P) = 1 \text{.}
\]
\end{theorem}

设 $\bar X$ 为补全 $X$ 的光滑射影曲线. 此公式是说 $\bar X$ 上的可逆层 $\sO(P)$ 的次数为 $1$.

\subsection{上同调方法}\label{IV:2-2}

\subsubsection{}\label{IV:2-2-1}

设 $X$ 为 (Noether) 概形. 回顾: 若 $X$ 的子概形 $Y$ 局部地 (在 $Y$ 上)
由 $X$ 中 $c$ 个方程组成的正则序列定义, 则称它是余维 $c$ 的局部完全交 (local complete intersection).
当 $X$ 为极大理想为 $\fm$ 的局部环 $A$ 的谱, $Y$ 的理想为 $\fa$ 时,
这意味着对 $i<\dim(\fa/\fm\fa)=c$, 有 $\Ext^i(A/\fa,A) = 0$,
而 $\fa$ 中在 $\fa/\fm\fa$ 内的像构成 $\fa/\fm\fa$ 的一组基的每个元素序列, 都是 $Y$ 的一个正则方程序列.

\subsubsection{}\label{IV:2-2-2}

设 $i:Y\hookrightarrow X$ 为余维 $c$ 的局部完全交.
我们要定义局部基本类 $\cl(Y)$, 它是带支集上同调层
$\eR^{2c}i^! \dZ/n(c)$ 的整体截面 (回顾 $\dZ/n(c)=\dmu_n^{\otimes c}$).

局部地, $Y$ 是一列 $c$ 个除子 $D_i$ 的交, 我们将 $\cl(Y)$ 定义为各 $\cl(D_i)$ 的杯积.
每个 $\cl(D_i)$ 的支集在 $D_i$ 中, 因而它们的乘积的支集在 $Y$ 中.
这一乘积局部地不依赖于 $D_i$ 的选取, 可由下面的 \ref{IV:2-2-3} 及下列不变性得到:
\begin{enumerate}[\indent a)]
  \item 与局部化相容;
  \item 不依赖于各 $D_i$ 的顺序 (各 $\cl(D_i)$ 的次数 $2$ 为偶数, 所以杯积可交换);
  \item 乘积仅依赖于“旗标”
    $D_1\supset D_1\cap D_2\supset \cdots\supset Y$.
\end{enumerate}
为证明 c), 注意存在乘积 (\ref{IV:1-2-1} 的一个变式)
\[
  \h_{D_1}^\bullet(X)\otimes \h_{D_1\cap D_2}^\bullet(D_1)\otimes \cdots \otimes \h_Y^\bullet(D_1\cap \cdots \cap D_{c-1}) \to \h_Y^\bullet(X) \text{;}
\]
各 $\cl(D_i)$ 的乘积也就是各 ($\cl(D_i)$ 限制到
$D_1\cap \cdots \cap D_{i-1})=(\cl(D_1\cap \cdots \cap D_i)$ 在
$D_1\cap \cdots \cap D_{i-1})$ 中的乘积.

\begin{lemma}\label{IV:2-2-3}
设 $A$ 为极大理想为 $\fm$ 的局部环,
$u=(u_1,\dots,u_c)$, $v=(v_1,\dots,v_c)$ 为生成同一理想 $\fa$ 的两个正则序列.
则存在由这样的序列构成的一列 $w_i$ ($1\leqslant i\leqslant N$) 将它们连接起来,
使 $w_{i+1}$ 可由 $w_i$ 经置换或只改变最后一个元素而得到.
\end{lemma}

既然允许置换, 则允许改变一个元素与允许只改变最后一个元素是等价的.
由 \ref{IV:2-2-1}, 我们只需验证: 在向量空间 $\fa/\fm\fa$ 中,
可通过一系列置换以及每次只改变一个向量的换基, 从一组基变到另一组基.
这确实成立, 因为线性群由对角矩阵、初等矩阵及置换矩阵生成.

\subsubsection{}\label{IV:2-2-4}

同样的方法使我们能够对每个局部可由 $c$ 个方程定义的 $Y\subset X$ 定义局部基本类
(在所需方程更少的地方, 此类为零).

\subsubsection{}\label{IV:2-2-5}

有半纯性结果时, 可以从这样的局部基本类
$\h^0\left(Y,\eR^{2 c} i^!\dZ/n(c)\right)$
过渡到整体基本类 $\h_Y^{2 c}\left(X,\dZ/n(c)\right)$ 中:

% 据法文原书印页141 (PDF146) 核对: (ii) 两处 i^! 均应为 k^!; 修复英文转录错误.
\begin{proposition}\label{IV:2-2-6}
\begin{enumerate}[(i)]
  \item 若对 $p<2 c$ 有 $\eR^p i^!\dZ/n=0$, 则
    \[
      \h_Y^{2 c}\left(X,\dZ/n(c)\right) \iso \h^0\left(Y,\eR^{2 c} i^! \dZ/n(c)\right) \text{.}
    \]
  \item 设 $Z$ 为 $Y$ 的闭部分, 在 $Y$ 中的补集为 $V$, $k$ 为 $Z$ 到 $X$ 的包含.
    若对 $p\leqslant 2 c$ 有 $\eR^p k^! \dZ/n = 0$, 则
    \[
      \h_Y^{2 c}\left(X,\dZ/n(c)\right) \hookrightarrow \h_V^{2 c}\left(X,\dZ/n(c)\right) \text{.}
    \]
    若对 $p\leqslant 2 c+1$ 有 $\eR^p k^!\dZ/n=0$, 则此箭头为同构.
\end{enumerate}
\end{proposition}

假设 $\eR^p i^!\dZ/n=0$ 等价于 $\eR^p i^!\dZ/n(c)=0$.
据此, (i) 由谱序列 $\h^p(Y,\eR^q i^!)\Rightarrow \h_Y^{p+q}(X,-)$ 得到.
若 $k_1$ 为 $Z$ 到 $Y$ 的包含, 则 $k=i k_1$, 从而
$\eR k^! \dZ/n=\eR k_1^! \eR i^!$,
而 $Z\subset Y$ 的上同调长正合列给出长正合列
\[\xymatrix{
  \cdots \ar[r]
    & \h^i\left(Z,\eR k^! \dZ/n(c)\right) \ar[r]
    & \h_Y^i\left(X,\dZ/n(c)\right) \ar[r]
    & \h_V^i\left(X,\dZ/n(c)\right) \ar[r]
    & \cdots
}\]
由此得到 (ii).

\subsubsection{推广}\label{IV:2-2-7}

对于任意有限型分离态射 $i$ 以及任意整数 $2 c$ (正或负),
只要将其中的 $\h_Y^{2 c}\left(X,\dZ/n(c)\right)$
换为 $\h^{2 c}\left(Y,\eR i^!\dZ/n(c)\right)$,
并对 $\h_V$ 作同样替换, 命题 \ref{IV:2-2-6} 仍然成立.

下列半纯性结果可由 \cite[1.8, 1.10, 1.15]{sga2} 得到. 我们回顾其证明.

\begin{theorem}\label{IV:2-2-8}
设有有限型 $S$-概形的态射
\[\xymatrix{
  Y \ar[r]^-i \ar[dr]_-f
    & X \ar[d]^-g \\
  & S
}\]
假定 $X$ 光滑, 纯相对维数为 $N$, 而 $Y$ 的各纤维维数 $\leqslant d$. 置 $c=N-d$.
\begin{enumerate}[\indent (i)]
  \item 对 $p<2 c$, 有 $\eR^p i^!\dZ/n=0$;
    将 $\dZ/n$ 换成层 $g^* \sF$ 也一样.
  \item 若在 $S$ 的稠密开子集 $U$ 上, $Y$ 的各纤维维数 $<d$,
    则对 $p\leqslant 2 c$ 有 $\eR^p i^!\dZ/n=0$.
    若此外 $U$ 的补集 $Z$ 局部地不使 $S$ 不连通,
    则对 $p\leqslant 2 c+1$ 有 $\eR^pi^!\dZ/n=0$.
\end{enumerate}
\end{theorem}

由于 $\eR g^! \sF=g^*\sF(N)[2N]$, 传递公式
$\eR i^!\eR g^!=\eR f^!$ 表明
% 法文原书印页142 (PDF147) 同样漏印右侧 =0; 此处补足, 见译注.
$\eR^{2 c+q} i^!(g^*\sF)=0\Leftrightarrow \eR^{-2 d+q}f^!(\sF)=0$.\footnote{译注: 原书印页142的等价式右端漏印 $=0$. 由 $\eR g^!\sF=g^*\sF(N)[2N]$ 及 $c=N-d$, 两侧所述正是相应上同调层的消失; 此处补足该等号.}
因此归结为研究 $f$, 即假定 $X=S$. 此时断言 (i) 就是 \cite[XVIII 3.17]{sga4}.

设 $Y'$ 为 $Z$ 在 $Y$ 中的逆像:
\[\xymatrix{
  Y' \ar[r]^-v \ar[d]^-f
    & Y \ar[d]^-f \\
  Z \ar[r]^-u
    & S
}\]
由 (i), 当 $p\leqslant -2 d+1$ 时, $\eR^p f^! \dZ/n$ 的支集在 $Y'$ 中;
谱序列 $\eR^a v^!\eR^b f^! \Rightarrow \eR^{a+b} (f v)^!$
表明它们与 $\eR^p(f v)^!\dZ/n=\eR^p(u f)^!\dZ/n$ 相同.
将 (i) 用于 $Y'/Z$ 以及谱序列
$\eR^a f^!\eR^b u^! \Rightarrow \eR^{a+b} (u f)^!$,
可知 (ii) 分别来自 $\eR^b u^! \dZ/n$ 在 $b=0$, 或在 $b=0$ 和 $1$ 时的消失.

\subsubsection{}\label{IV:2-2-9}

Grothendieck 猜想 \ref{IV:2-2-8} 有如下绝对对应形式
(半纯性猜想, 是纯性猜想的推论): 若 $Y$ 在正则 $X$ 中的余维 $\geqslant c$,
则对 $i<2 c$ 有 $\h_Y^i(X,\dZ/n)=0$, 至少在 $X$ 优良时成立.

\subsubsection{}\label{IV:2-2-10}

现在可以定义域 $k$ 上光滑的 $X$ 中余维 $c$ 的闭链 $Y$ 的类.
写成 $Y=\sum d_i Y_i$, 其中各 $Y_i$ 约化且不可约.
$Y_i$ 的某个开子集 $U_i$ 是 $X$ 中的局部完全交, 且其补集的余维 $>c$.
于是可以定义 $U_i$ 的局部基本类.
由 \ref{IV:2-2-6} 和 \ref{IV:2-2-8}, 它来自唯一的基本类
$\cl(Y_i)\in \h_{Y_i}^{2 c}\left(X,\dZ/n(c)\right)$, 我们置
\[
  \cl(Y) = \sum d_i \cl(Y_i) \in \h_{|Y|}^{2 c}\left(X,\dZ/n(c)\right) \text{.}
\]

\subsubsection{}\label{IV:2-2-11}

在复解析几何中, Baum、Fulton 和 Mac Pherson \cite{bfm75} 的一个构造,
可以不加限制地定义局部完全交 $Y\subset X$ 的上同调类.
假定 $Y$ 的纯余维为 $c$, 以 $\sN$ 表示 $Y$ 上的向量丛, 即 $Y$ 在 $X$ 中的法丛.
其局部截面为 $(\sI/\sI^2)^\vee$ 的截面, 其中 $\sI$ 是 $Y$ 的理想层.
回顾: $X$ 上 $C^\infty$ 函数层的定义, 局部地通过 $X$ 到 $\dC^n$ 的局部嵌入给出:
将 $\dC^n$ 上的 $C^\infty$ 函数层模去由定义 $X$ 的方程的实部和虚部生成的理想,
再将所得商限制到 $X$.
% 法文印页143 (PDF148) 中 C^infty 表示光滑性, N 为丛名; 分开数学文字片段, 不改符号.
丛 $\sN$ 可延拓为 $Y$ 在 $X$ 中某邻域 $U$ 上的复 $C^\infty$ 向量丛 $N$;
且 $U$ 充分小时, 存在 $N$ 的截面 $f$, 其零点集为 $Y$,
并使在 $Y$ 上 $d f:\sN\to N$ 为恒等映射.
$N$ 和 $f$ 的两种选取在 $Y$ 的充分小邻域上同伦.

$N$ 的截面 $0$ (记为 $z:U\to N$) 的上同调类 $\cl(U)$ 已有定义:
$U$ 是 $N$ 中的局部完全交, 且在 $i\ne 2 c$ 时有 $\eR^i z^!\dZ = 0$,
又有 $\eR^{2 c} z^!\dZ = \dZ$.
我们有 $f^*:\h_U^\bullet(N,\dZ) \to \h_Y^\bullet(X,\dZ)$, 并置
\[
  \cl(Y) = f^*\cl(U) \text{.}
\]

只要 $X$ 的解析子空间 $Y$ 带有如下法向结构, 同一构造仍适用:
$Y_\text{red}$ 上秩 $c$ 的局部自由层 $\sC$, 以及满态射
$\sC\to \sI/\sI\cdot \sI_\text{red}$.

\subsection{同调方法}\label{IV:2-3}

\subsubsection{}\label{IV:2-3-1}

设 $f:Y\to S$ 为有限型平坦态射, 纤维维数 $\leqslant d$.
在 \cite[XVIII 2.9]{sga4} 中, 我们定义了迹态射
\begin{equation*}\tag{2.3.1.1}\label{IV:eq:2-3-1-1}
  \tr_f:\eR^{2 d} f_! \dZ/n(d) \to \dZ/n \text{.}
\end{equation*}
在 $i>2 d$ 时有 $\eR^i f_!\dZ/n(d)=0$,
在 $i<-2 d$ 时有 $\eR^i f^!\dZ/n=0$.
结合 $\eR f_!$ 与 $\eR f^!$ 之间的伴随, 得到同构
\begin{equation*}\tag{2.3.1.2}\label{IV:eq:2-3-1-2}
\begin{aligned}
  \hom\left(\eR^{2 d}f_!\dZ/n(d),\dZ/n\right)
    &= \hom\left(\eR f_!\dZ/n(d),\dZ/n[-2 d]\right) \\
    &= \hom\left(\dZ/n,\eR f^!\dZ/n(-d)[-2 d]\right) \\
    &= \h^0\left(Y,\eR^{-2 d}f^!\dZ/n\right) \text{.}
\end{aligned}
\end{equation*}
至少当 $S$ 为一点时, $\tr_f$ 在最后两个群中的像适合称为 $Y$ 的基本类 (在同调中).\footnote{译注: 原书印页144的末项漏去扭转 $(-d)$; 由前一项取零次上同调, 应读作 $\h^0(Y,\eR^{-2d}f^!\dZ/n(-d))$. 此处保留原式.}

我们仍以 $\tr_f$ 表示 $\tr_f$ 在第二个群中的像,
以及通过函子性由其得到的态射. 例如, 当 $Y/S$ 紧合时, 有态射
\begin{equation*}\tag{2.3.1.3}\label{IV:eq:2-3-1-3}
\xymatrix{
  \h^i\left(Y,\dZ/n(d)\right) = \h^i\left(S,\eR f_!\dZ/n(d)\right) \ar[r]
    & \h^{i-2 d}(S,\dZ/n) \text{.}
}
\end{equation*}

假定 $Y$ 包含在 $S$ 上光滑且纯相对维数为 $N$ 的 $X$ 中, 并置 $c=N-d$.
\[\xymatrix{
  Y \ar@{^{(}->}[r]^-i \ar[dr]_-f
    & X \ar[d] \\
  & S
}\]
有 $\eR f^! = \eR i^!\eR g^!$, 以及
$\eR g^! \dZ/n=\dZ/n(-N)[-2 N]$.\footnote{译注: 此负扭转、负移位为原书印页144的笔误. 按 \ref{IV:2-2-8} 的光滑态射约定, 应为 $\eR g^!\dZ/n=\dZ/n(N)[2N]$; 结合前式的 $(-d)[-2d]$, 才得到下式的 $(c)[2c]$.}
$Y$ 的基本类与
\[
  \h^0\left(Y,\eR i^!\dZ/n(c)[2 c]\right)=\h_Y^{2 c}\left(X,\dZ/n(c)\right)
\]
中的一个元素等同, 即 $Y$ 在 $X$ 中的类 $\cl(Y)$.
下面将看到, 它只依赖于 $Y\subset X$, 而不依赖于 $X$ 到 $S$ 的投影;
且当 $Y$ 为局部完全交时, 它诱导上一小节的局部类.

现在明确写出将 $\cl(Y)$ 映为 $\tr_f$ 的同构
\[
  \h_Y^{2 c}\left(X,\dZ/n(c)\right) \iso \hom\left(\eR^{2 d}f_!\dZ/n(d),\dZ/n\right)
\]
: 通过同构
\begin{align*}
  \h_Y^{2 c}\left(X,\dZ/n(c)\right)
    &= \h_Y^{2 c}\left(\eR i^!\dZ/n(c)\right) \\
    &= \hom_Y\left(\dZ/n(d)[2 d],\eR i^!\dZ/n(N)[2 N]\right) \\
    &= \hom_X\left(\dZ/n(d)[2 d]_Y,\dZ/n(N)[2 N]\right) \text{,}
\end{align*}
它就是下图的上行:
\[\xymatrix@C=1.8em{
  \h_Y^{2 c} \ar[r] \ar[dr]_-{(3)}
    & \hom\left(\eR f_!\dZ/n(d)[2 d],\eR f_! \eR i^!\dZ/n(N)[2 N]\right) \ar[r]^-{(1)} \ar[d]^-{\tr_i} \ar@{}[dr]|-{(2)}
    & \hom\left(\eR f_!\dZ/n(d)[2d],\dZ/n\right) \ar@{=}[d] \\
  & \hom\left(\eR g_!\dZ/n(d)[2 d]_Y,\eR g_!\dZ/n(N)[2 N]\right) \ar[r]^-{\tr_g}
    & \hom\left(\eR g_!\dZ/n(d)[2 d]_Y,\dZ/n\right) \text{,}
}\]
其中 (1) 为伴随箭头 $\eR f_!\eR f^!\to \text{id}$.
(2) 的交换性表达同构 $\eR f^!=\eR i^!\eR g^!$ 是通过伴随定义的.
最后, 由 (3) 得到的箭头
\[
  \h_Y^{2c}\left(X,\dZ/n(c)\right) \otimes \eR^{2 d} f_!\dZ/n(d) \to \eR^{2 N} g_!\dZ/n(N)
\]
解释为杯积 (参见 \ref{IV:1-2}).

\begin{definition}\label{IV:2-3-2}
类 $\cl(Y)\in\h_Y^{2 c}\left(X,\dZ/n(c)\right)$ 由下列性质刻画:
对于 $\eR^{2 d}f_! \dZ/n(d)$ 的每个局部截面 $u$, 有
\[
  \tr_f(u) = \tr_g\left(\cl(Y)\smallsmile u\right) \text{.}
\]
\end{definition}

由于 $\cl(Y)$ 在
$\hom\left(\eR f_!\dZ/n(d)[2 d],\dZ/n\right)$ 中已有像 $\tr_f$,
公式 \ref{IV:2-3-2} 对通过函子性由 $\tr_f$ 得到的箭头也成立.
例如, 当 $X$ 在 $S$ 上紧合且 $u\in \h^\bullet(Y,\dZ/n)$ 时,
该公式在 $\h^\bullet\left(S,\dZ/n(-d)\right)$ 中成立.
对于 $v\in \h^\bullet(X,\dZ/n)$, 它给出
\[
  \tr_f(i^* v) = \tr_g\left(\cl(Y)\smallsmile v\right) \text{,}
\]
其中杯积可在 $\h^\bullet(X,\dZ/n)$ 中计算.

\subsubsection{}\label{IV:2-3-3}

我们将在更一般的假设下定义迹态射, 因而也定义类 $\cl(Y)$. 考虑图
\[\xymatrix{
  Y \ar@{^{(}->}[r] \ar[dr]_-f
    & X \ar[d]^-g \\
  & S
}\]
其中 $g$ 光滑且纯相对维数为 $N$, $Y$ 在 $X$ 中闭, 各纤维维数 $\leqslant d$.
再给 $Y$ 配备如下类型的“权” $\sK$:
$\sK\in\eD_{\text{parf},X}^b$ 是 $X$ 上 $\sO$-模层的有界复形,
在 $S$ 上 (或在 $X$ 上, 两者等价) 具有有限 $\tor$-维数,
其上同调层凝聚且支集在 $Y$ 中.
我们要定义态射 $\tr_{f,\sK}:\eR^{2 d}f_! \dZ/n(d) \to \dZ/n$.
当 $Y$ 在 $S$ 上平坦且 $\sK=\sO_Y$ 时, 它就是前面的迹态射.
一般地, 它只依赖于 $\sK$ 在 $Y$ 的各不可约分支的一般点 $y$ 处的长度
($\operatorname{lg}_y(\sK) = \sum (-1)^i \operatorname{lg} \h^i(\sK)_y$).
% 核对记录: 英文源在此注明上同调符号是否应为花体 H 尚不明确; 保留原式.
最后, 以 $\cl(Y,\sK)$ 表示支集在 $Y$ 中且满足
\begin{equation*}\tag{2.3.3.1}\label{IV:eq:2-3-3-1}
  \tr_g\left(\cl(Y,\sK)\smallsmile u\right) = \tr_{f,\sK}(u) \text{.}
\end{equation*}
的类.

\subsubsection{}\label{IV:2-3-4}

$\tr_{f,\sK}$ 的构造与 \cite[XVIII.2]{sga4} 中的构造平行; 我们只指出其主要步骤.

\paragraph{A}
$d=0$ ($f$ 拟有限). 设 $x$ 为 $Y$ 的几何点, $s$ 为其在 $S$ 中的像,
$\sK_{(x)}$ 为 $\sK$ 在 $X$ 于 $x$ 的严格局部化上的逆像,
$\sK_{(x)s} = \sK_{(x)}\lotimes_{\sO_{S,s}^h} k(s)$ 为其在 $s$ 处几何纤维上的逆像,
并置 $n(x)=\sum (-1)^i \dim_{k(s)} \h^i(\sK_{(x)s})$.
函数 $x\mapsto n(x)$ 是 $f$ 的一个加权, 我们取对应的迹态射 \cite[XVII.6.2.5]{sga4}.
加权 $n(x)$ 及 $\tr_{f,\sK}$ 的形成与任意基变换相容.

\paragraph{B}
一般情形. 若 $u:X\to \dA_S^d$ 使 $u i$ 拟有限,
则置 $\tr_{f,\sK} = \tr_{\dA_S^d}\circ \tr_{u i,\sK}$.
此迹态射显然与每个基变换 $S'/S$ 相容.
所以为证明它不依赖于 $u$, 可假定 $S$ 为代数闭域 $k$ 的谱.
此时, 设 $Y_i$ 为 $Y$ 的不可约分支, $\operatorname{lg}_i$ 为 $\sK$ 在 $Y_i$ 一般点处的长度.
若 $\alpha_i$ 为 $(Y_i)_\text{red}$ 到 $Y$ 的包含, 则有
% 法文印页147 (PDF152) 的公式号为 (2.3.4.1), 英文漏去末尾 .1.
\begin{equation*}\tag{2.3.4.1}\label{IV:eq:2-3-4}
  \tr_{f,\sK} = \sum \operatorname{lg}_i \tr_{Y_{i,\text{red}/S}} \alpha_i^\ast
\end{equation*}
(此公式由 $\tr_{u i}$ 的对应的容易验证的公式得到).

这样在 $Y$ 上局部地定义了 $\tr_{f,\sK}$; 随后按 \cite[XVIII.2.9]{sga4} 继续即可.

\begin{lemma}\label{IV:2-3-5}
若 $g=g'g'':X\xrightarrow{g''} S'\xrightarrow{g'} S$,
其中 $g'$ 和 $g''$ 光滑, 纯相对维数分别为 $N'$ 和 $N''$,
且 $Y$ 在 $S'$ 上的各纤维余维仍 $\geqslant c$,
则类 $\operatorname{cl}(Y,\sK)$ 无论相对于 $g$ 还是相对于 $g''$ 计算都相同.
\end{lemma}

设 $f'=g'' i$, $d'=d-N'$.
公式 $\tr_g=\tr_{g'}\tr_{g''}$ (\cite[XVIII.2.9]{sga4} 变式 3) 保证下图交换:
% 据法文印页147 (PDF152) 补回左上因子 f'_! 的撇号.
\[\xymatrix{
  \h_Y^{2 c}\left(X,\dZ/n(c)\right) \ar[d] \\
  \hom_{S'}\left(\eR f_!'\dZ/n(d')[2 d'],\eR g_!''\dZ/n(N'')[2 N'']\right) \ar[r] \ar[d]^-{\eR g_!'(N')}
    & \hom_{S'}\left(\eR^{2 d'} f_!' \dZ/n,\dZ/n\right) \ar[d]^-{\tr_{g'}\circ \eR^{2 N'} g_!'} \\
  \hom_S\left(\eR f_!\dZ/n(d)[2d],\eR g_!'\dZ/n(N)[2 N]\right) \ar[r]
    & \hom_S\left(\eR^{2 d} f_!\dZ/n,\dZ/n\right)
}\]
\footnote{译注: 图中另有原书印页147的记号遗漏: 右列的 $\eR^{2d'}f'_!\dZ/n$ 与 $\eR^{2d}f_!\dZ/n$ 分别须带扭转 $(d')$ 与 $(d)$, 才能与左列的迹态射对应; 左下项的第二因子应为 $\eR g_!\dZ/n(N)[2N]$, 因 $g$ 的定义域为 $X$, 而 $g'$ 的定义域为 $S'$. 此处保留这些原式.}
而对应的、容易验证的公式 $\tr_{f,\sK}=\tr_{g'}\tr_{f',\sK}$
保证 $\tr_{f',\sK}$ 的像是 $\tr_{f,\sK}$. 断言随之成立.

\begin{lemma}\label{IV:2-3-6}
若 $Y$ 是 $X$ 中的除子, 则 \ref{IV:2-3-2} 的类与 \ref{IV:2-1-2} 的类相同.
\end{lemma}

由半纯性 (\ref{IV:2-2-6}(i) 和 \ref{IV:2-2-8}(i)), 问题在 $Y$ 上是局部的;
再用 \ref{IV:2-3-5} 将 $S$ 换成适当的 $S'$, 就可假定 $N=1$.
再次利用半纯性 (\ref{IV:2-2-6}(ii) 和 \ref{IV:2-2-8}(ii)),
只需在 $S$ 的一般点上证明 \ref{IV:2-3-6}.
再经平展局部化, 可归结到 $S$ 为代数闭域 $k$ 的谱的情形.
由于 \ref{IV:2-3-2} 和 \ref{IV:2-1-2} 的类都对 $Y$ 可加,
可归结到 $Y$ 为 $k$ 上光滑曲线的闭点的情形.
此时, 刻画性质 \ref{IV:2-3-2} 成立这一事实, 就是基本相容性 \ref{IV:2-1-5}.

\begin{lemma}\label{IV:2-3-7}
$\cl(Y,\sK)$ 的形成与任意基变换 $S'/S$ 相容.
\end{lemma}

这由迹态射的同一断言得到.

\begin{theorem}\label{IV:2-3-8}
(i) $\cl(Y,\sK)$ 仅依赖于 $Y\subset X$ 及 $\sK$ 在 $Y$ 的一般点处的长度.
此类对 $\sK$ 可加. 当 $\sK=\sO_Y$ 且 $Y$ 为局部完全交时,
它诱导 \ref{IV:2-2-2} 的局部类.

(ii) 设有交换图
\[\xymatrix{
  X' \ar[r]^-u \ar[d]
    & X \ar[d] \\
  S' \ar[r]
    & S
}\]
其中 $X'$ 在 $S'$ 上光滑. 若 $u^{-1}(Y)$ 的各纤维余维仍 $\geqslant c$, 则
\[
  \cl\left(u^{-1}(Y),\eL u^\ast\sK\right) = u^\ast \cl(Y,\sK) \text{.}
\]

(iii) 设 $Y'\subset X$ 的各纤维余维 $\geqslant c'$,
$\sK'$ 为 $Y'$ 上的权, 并假定 $Y\cap Y'$ 的各纤维余维 $\geqslant c+c'$. 则
\[
\cl\left(Y\cap Y',\sK\lotimes \sK'\right) = \cl(Y,\sK) \smallsmile \cl(Y',\sK') \text{.}
\]
\end{theorem}

% 据法文印页148--149 (PDF153--154) 恢复英文以 A/B/C 替换的原段落号.
\paragraph{(2.3.8.1)}
给定上述 $Y\subset X/S$, 由半纯性 (\ref{IV:2-2-6}, \ref{IV:2-2-8}) 可知,
要验证 $\h_Y^{2 c}\left(X,\dZ/n(c)\right)$ 中的两个类相同,
只需在 $Y$ 的各一般点处局部地验证;
甚至只需在每次基变换 $s\to S$ ($s$ 为 $S$ 的几何一般点) 后, 在 $Y$ 的各一般点处验证.

\paragraph{(2.3.8.2) 当 $S=S'$ 且 $u$ 为闭浸入时 (ii) 的证明}
由局部化 (2.3.8.1), 可假定 $X'$ 是光滑态射 $v:X\to \dA_S^{N'}$ 下零截面 $0$ 的逆像:
\[\xymatrix{
  X' \ar[r]^-u \ar[d]
    & X \ar[d]^-v \\
  S \ar[r]^-0
    & \dA_S^{N'} \ar[r]^-\pi
    & S
}\]
对 $g=\pi v$ 应用 \ref{IV:2-3-5}, 对基变换 $S\to \dA_S^{N'}$ 应用 \ref{IV:2-3-7}.

\paragraph{(2.3.8.3) 在乘积情形中:}
$X=X'\times_S X''$, $Y=Y'\times_S Y''$, $\sK=\sK'\lotimes_S \sK''$, 有
$\cl(Y,\sK) = \cl(Y',\sK')\smallsmile \cl(Y'',\sK'')$
(外杯积, 即 $\operatorname{pr}_1^\ast \cl(Y',\sK')\smallsmile \operatorname{pr}_2^\ast \cl(Y'',\sK'')$).

$\tr_g$ 的函子性质 \cite[XVIII.2.12]{sga4} 蕴涵下图交换:
\small
\[\xymatrix@=0.5cm{
  \h_{X'}^{2c'}\otimes \h_{X''}^{2 c''}\ar[r] \ar[d]
    & {\begin{gathered}\hom_S\left(\eR f_!' \dZ/n(d')[2 d'],\eR g_!'\dZ/n(N')[2 N']\right)\\ {}\otimes \cdots\end{gathered}} \ar[r] \ar[d]
    & {\begin{gathered}\hom_S\left(\eR^{2 d'} f_!' \dZ/n,\dZ/n\right)\\ {}\otimes \cdots\end{gathered}} \ar[d] \\
  \h_X^{2 c} \ar[r]
    & \hom_S\left(\eR f_!\dZ/n(d)[2 d],\eR g_!\dZ/n(N)[2 N]\right) \ar[r]
    & \hom_S\left(\eR^{2 d}f_!\dZ/n,\dZ/n\right)
}\]
\normalsize
并可验证 $\tr_f$ 是 $\tr_{f'}\otimes \tr_{f''}$ 的像, 断言由此得到.

\paragraph{(ii) 的证明}
由 (\ref{IV:2-3-7}) 并先作一次基变换, 可假定 $S=S'$.
将 $u$ 通过其图像分解:
\[\xymatrix{
  X' \ar[r]^-{(\text{id},u)}
    & X'\times X \ar[r]^-{\operatorname{pr}_2}
    & X \text{.}
}\]
于是分别处理 $\operatorname{pr}_2$ 与 $(\text{id},u)$:
前者归于 (2.3.8.3), 取 $Y'=X'$, $\sK'=\sO_{X'}$
(此时 $\cl(Y')=1\in \h^0(X',\dZ/n)$); 后者归于 (2.3.8.2).

\paragraph{(iii) 的证明}
设 $\delta:X\to X\times_S X$ 为对角包含.
利用公式 $\delta^\ast(Y_1\times_S Y_2) = Y_1\cap Y_2$, 归结到 (2.3.8.2) 和 (2.3.8.3).

\paragraph{(i) 的证明}
基变换 $S_\text{red}\to S$ 将 $X$ 换为 $X_\text{red}$;
从而可假定 $X$ 和 $S$ 约化.
然后在 $Y$ 的一般点附近局部化, 可归结到 $Y_\text{red}$ 不可约且为 $c$ 个除子 $D_i$ 的完全交的情形.
同样的局部化, 结合拟有限情形中 $\tr$ 的定义, 表明
$\cl(Y,\sK)=\operatorname{lg}\cl(Y_\text{red},\sO_{Y_\text{red}})$,
其中 $\operatorname{lg}$ 为 $\sK$ 在 $Y$ 一般点处的长度.
因此由 (iii), 有
\[
  \cl(Y,\sK) = \operatorname{lg} \prod_i \cl(D_i,\sO_{D_i})
\]
结论即由 \ref{IV:2-3-6} 得到.

\subsubsection{注}\label{IV:2-3-9}

当 $S$ 为域的谱、闭链的类按 \ref{IV:2-2-10} 定义时,
断言 (iii) 是说, 若两个闭链 $Y'$ 和 $Y''$ 按正确的维数相交, 则
\[
  \cl(Y'\cap Y'') = \cl(Y')\smallsmile \cl(Y'')\text{,}
\]
其中 $Y'\cap Y''$ 的各分支的重数须以 $\tor$ 的交错和计算.

\subsubsection{注}\label{IV:2-3-10}

设 $X/\spec(k)$, $k$ 代数闭. 若 $X$ 中两个余维 $c$ 的闭链代数等价,
则上述相对理论的存在表明, 它们的类在 $\h^{2 c}(X,\dZ/n(c))$ 中的像,
是某个适当选取的连通 $S$ 上的常值层 $\h^{2 c}(X,\dZ/n(c))$
(即 $f=\operatorname{pr}_2:X\times S\to S$ 的层 $\eR^{2 c} f_\ast \dZ/n(c)$)
在 $S$ 上的一个截面于两个点处的纤维:
两个代数等价的闭链在 $\h^{2 c}(X,\dZ/n(c))$ 中有相同的类.

\section{应用: 紧合光滑情形的 Lefschetz 迹公式}\label{IV:3}

\subsection{}\label{IV:3-1}

设 $X$ 和 $Y$ 为代数闭域 $k$ 上的紧合光滑代数簇, 纯维数分别为 $N$ 和 $M$.
通过 K\"unneth 同构 $\h^\bullet(X\times Y)=\h^\bullet(X)\otimes \h^\bullet(Y)$,
迹态射 $\h^\bullet(X\times Y)(M)\to \h^\bullet(X)$ 就是
$\h^\bullet(X)\otimes ($迹态射 $\h^\bullet(Y)(M)\to \dQ_\ell)$.
由于此态射是偶数次齐次的, 没有符号问题. 记为 $\int_Y$.

\subsection{}\label{IV:3-2}

类 $\eta\in \h^{2 q}(X\times Y)(q)$ 定义一个次数为 $2(q-M)$ 的态射
$\eta_{XY}:\h^\bullet(Y)(M-q)\to \h^\bullet(X)$, 其公式为
\[
  \beta \mapsto \int_Y \eta\cdot \operatorname{pr}_2^\ast \beta \text{.}
\]

\begin{proposition_}\label{IV:3-3}
设 $p$ 和 $q$ 为两个整数, $p+q=N+M$, 并设
$\varepsilon\in \h^{2 p}(X\times Y)(p)$,
$\eta\in \h^{2 q}(X\times Y)(q)$.
将 $\eta_{X Y}\varepsilon_{Y X}:\h^\bullet(X)\to \h^\bullet(X)$ 的迹理解为 \ref{IV:1-3} 的意义, 则有
\[
  \tr_{X\times Y}(\eta\cdot \varepsilon) = \tr\left(\eta_{X Y}\varepsilon_{Y X},\h^\bullet(X)\right) \text{.}
\]
\end{proposition_}

为消除符号的困扰, 只保留 $\h^\bullet$ 分次的底层 $\dZ/2$-分次会更方便.
为免不断写出 Tate 扭转, 再固定同构 $\dQ_\ell(1)=\dQ_\ell$.
设 $\alpha_Y$ 为使下图交换的同态 $\h^\bullet(Y)\to \h^\bullet(Y)^\vee$ (它是同构):
\[\xymatrix{
  \h^\bullet(Y)\otimes \h^\bullet(Y) \ar[r] \ar[d]^-{\alpha_Y\otimes 1}
    & \h^\bullet(Y) \ar[r]^-{\tr}
    & \dQ_\ell \ar@{=}[d] \\
  \h^\bullet(Y)^\vee \otimes \h^\bullet(Y) \ar[rr]^-{\text{(\ref{IV:eq:1-3-1})}}
    & & \dQ_\ell \text{.}
}\]
映射 $\eta_{X Y}$ 为 $\eta$ 在
\[\xymatrix{
  \h^\bullet(X\times Y) = \h^\bullet(X)\otimes \h^\bullet(Y) \ar[r]^-{\alpha_Y}
    & \h^\bullet(X) \otimes \h^\bullet(Y)^\vee \ar@{=}[r]^-{\text{(\ref{IV:eq:1-3-2})}}
    & \hom\left(\h^\bullet(Y),\h^\bullet(X)\right) \text{,}
}\]
下的像, $\varepsilon$ 也一样.
于是迹的定义 \ref{IV:eq:1-3-4} 将 \ref{IV:3-3} 归结为公式 $\tr_{X\times Y}=\tr_X\otimes \tr_Y$.

\subsection{注}\label{IV:3-4}

\ref{IV:3-3} 的证明没有用到 $\alpha_X$ 与 $\alpha_Y$ 为同构这一事实.
不假定 $X$ 和 $Y$ 光滑, 它仍成立.
不过在非光滑情形, 难以找到值得应用 \ref{IV:3-3} 的\emph{上同调}类.

\subsection{注}\label{IV:3-5}

若 $\varepsilon,\eta$ 是 $X\times Y$ 上余维分别为 $p$ 和 $q$ 的代数闭链 $A$ 和 $B$ 的上同调类,
且 $A\cap B$ 的维数为 $0$, 则 $\tr_{X\times Y}(\eta\varepsilon)$
就是以 $\tor$ 的交错和计算的相交重数 $A\cdot B$ (\ref{IV:2-3-9}).

\begin{proposition_}\label{IV:3-6}
当 $q=M$, $\eta=\cl(B)$ 为在 $X$ 上有限的子概形 $B\subset X\times Y$ 的类时,
映射 $\eta_{X Y}$ 是复合
\[\xymatrix{
  \eta_{X Y} : \h^\bullet(Y) \ar[r]^-{\operatorname{pr}_2^\ast}
    & \h^\bullet(B) \ar[r]^-{\tr_{B/X}}
    & \h^\bullet(X) \text{.}
}\]
\end{proposition_}

事实上, 由 \ref{IV:3-2} 及 $\sK=\sO_B$ 时的 \eqref{IV:eq:2-3-3-1}, 有
\[
  \eta_{X Y}(\beta) = \tr_{X\times Y/X}\left(\cl(B)\cdot \operatorname{pr}_2^\ast\beta\right) = \tr_{B/X}\left(\operatorname{pr}_2^\ast \beta\right) \text{.}
\]

最重要的情形是 $B$ 为映射 $f:X\to Y$ 的图像.
此时 \ref{IV:3-6} 表明 $\eta_{X Y}=f^\ast$.

\begin{corollary_}\label{IV:3-7}
设 $f$ 为代数闭域 $k$ 上紧合光滑的 $X$ 的自同态.
假定 $f$ 的不动点孤立. 则迹 $\sum (-1)^i \tr(f^\ast,\h^i(X))$
是 $f$ 的不动点个数, 每点按其重数计数.
\end{corollary_}

在 \ref{IV:3-5} 中取 $A=$ 恒等映射的图像, $B=$ $f$ 的图像,
并应用 \ref{IV:3-3}、\ref{IV:3-6}.

\begin{corollary_}\label{IV:3-8}
设 $X$ 为代数闭域 $k$ 上的曲线, 由 $\dF_q$ 上的 $X_0$ 经标量扩张得到,
$f$ 为 Frobenius 态射. 则 $\sum (-1)^i \tr\left(f^\ast,\h_c^i(X)\right)$
是 $f$ 的不动点个数.
\end{corollary_}

可以将 $X$ 换为 $X_\text{red}$, 从而假定 $X$ 约化.
设 $U$ 为 $X$ 的开子集, 在其上 $X$ 光滑且维数为 $1$, $\bar U$ 为 $U$ 的光滑射影补全.
$(X,U)$ 与 $(\bar U,U)$ 的上同调长正合列给出
\[
  \tr\left(f^\ast,\h_c^\bullet(X)\right) = \tr\left(f^\ast,\h_c^\bullet(U)\right) + \tr\left(f^\ast,\h_c^\bullet(X\setminus U)\right)
\]
以及
\[
  \tr\left(f^\ast,\h_c^\bullet(\bar U)\right) = \tr\left(f^\ast,\h_c^\bullet(U)\right) + \tr\left(f^\ast,\h_c^\bullet(\bar U\setminus U)\right) \text{.}
\]
不动点个数也满足相同的公式.
由于公式 \ref{IV:3-8} 对维数 $0$ 的概形显然成立, 只需对 $\bar U$ 证明 \ref{IV:3-8}.
这一情形由 \ref{IV:3-7} 得到. $d f=0$ 这一事实保证所有不动点的重数为一.

\subsection{注}\label{IV:3-9}

当 $X$ 在代数闭域 $k$ 上的维数 $\leqslant 1$, 且 $f:X\to X$ 满足 $d f=0$ 时,
不难验证 \ref{IV:3-8} 仍成立.
